% META: {"id": "TCOMPL-BAYES-EX-014", "chapitre": "TCOMPL-INFERENCE-BAYESIENNE", "type_objet": "exercice", "capacites": ["TCOMPL-BAYES-C4"], "capacites_codes": ["C4"], "methodes": ["M2","M5"], "parcours": 1, "competences": ["raisonner","calculer"], "duree_min": 18, "mode_creation": "ex_nihilo", "sources_inspiration": [], "fichier_tex": "chapitres/TCOMPL-INFERENCE-BAYESIENNE/exercices/TCOMPL-BAYES-EX-014.tex", "corrige_tex": "chapitres/TCOMPL-INFERENCE-BAYESIENNE/corriges/TCOMPL-BAYES-CO-014.tex", "status": "approved"}
% BEGIN-VERIFY
% from sympy import *
% x = symbols('x')
% f = x**2 * exp(-x)
% fp = diff(f, x)
% assert simplify(fp) == (2 - x)*x*exp(-x)
% assert solve(fp, x) == [0, 2]
% assert f.subs(x, 2) == 4*exp(-2)
% assert limit(f, x, oo) == 0
% assert limit(f, x, -oo) == oo
% END-VERIFY
\begin{exercice}{TCOMPL-BAYES-EX-014}{2}{18}
Soit $f(x) = x^2\,\mathrm{e}^{-x}$ definie sur $\mathbb{R}$.
\begin{enumerate}
  \item Calculer $f'(x)$ et etudier son signe.
  \item Determiner les limites de $f$ en $\pm\infty$.
  \item Dresser le tableau de variations de $f$.
\end{enumerate}

\end{exercice}
